\documentclass[11pt,a4paper]{article}
\usepackage[margin=1in]{geometry}
\usepackage[T1]{fontenc}
\usepackage[utf8]{inputenc}
\usepackage{amsmath,amssymb,bm}
\usepackage{booktabs}
\usepackage{graphicx}
\usepackage{microtype}
\usepackage[font=small,labelfont=bf]{caption}
\usepackage{xcolor}
\usepackage{url}
\usepackage[colorlinks=true,linkcolor=blue!50!black,urlcolor=blue!50!black]{hyperref}

\newcommand{\R}{\mathbf{R}}
\newcommand{\A}{\mathbf{A}}
\newcommand{\G}{\mathbf{G}}
\newcommand{\Cd}{\mathbf{C}_d}
\newcommand{\Cm}{\mathbf{C}_m}
\newcommand{\Lop}{\mathbf{L}}
\newcommand{\M}{\mathbf{M}}
\newcommand{\W}{\mathbf{W}}
\newcommand{\bfm}{\bm{m}}
\newcommand{\bfd}{\bm{d}}
\newcommand{\bfs}{\bm{s}}
\newcommand{\bfu}{\bm{u}}
\newcommand{\note}[1]{\textcolor{red!60!black}{\textbf{#1}}}

\title{\bf Resolution-driven slip inversion from InSAR line-of-sight
displacement\\[2mm] \large Method, mathematics, and findings from a six-scene
Venezuela study}
\author{\texttt{nisar\_tools/slip}}
\date{30 August 2026}

\begin{document}
\maketitle

\begin{abstract}
\noindent
This note sets out the mathematics behind the slip-inversion module in
\texttt{nisar\_tools} --- a Python port of, and extension to,
\href{https://github.com/x3zou/SlipSolve-Curve}{SlipSolve-Curve} (x3zou, MIT) ---
and records what has been measured with it. The forward
problem is standard: line-of-sight displacement from dislocation sources on a
triangulated fault. Two pieces have no published equivalent that we could find.
The first is a \emph{per-(triangle, observation) adaptive quadrature} rule for
layered elastic Green's functions, which makes a layered engine cost about the
same as a closed-form homogeneous half-space and therefore removes the last
computational argument for assuming a uniform crust. The second is a
\emph{joint loop} in which a single model-resolution matrix drives both the mesh
refinement and the data subsampling, alternating --- the published arrangement
keeps these two sequential and separate.

The findings are more equivocal than the methods. Measured against matched
effective rank, the mesh design, the sampling schedule and the data covariance
all turn out to be second-order: the regularisation weight $\lambda$ dominates
the answer, four independent ways. We report that as a result rather than
bury it, and quantify it: at the chosen operating point a typical slip parameter
is roughly 97\% regulariser, and the regularisation bias exceeds the propagated
data variance by an order of magnitude at every usable $\lambda$. Scored against
planted truth under measured noise, the L-curve corner that set that operating
point sits at about a quarter of the effective rank the data supports, and costs
40\% in model error.
\end{abstract}

\tableofcontents

\section{Scope and notation}

We solve for slip on a fault surface from surface displacement measured by
InSAR. Throughout, $\bfm$ is the parameter vector, $\bfd$ the data, and $\A$
the design matrix that maps one to the other. Vectors are columns; $'$ denotes
transpose. Table~\ref{tab:symbols} collects the symbols.

\begin{table}[htbp]
\centering
\small
\begin{tabular}{ll}
\toprule
symbol & meaning \\
\midrule
$\bfd \in \mathbb{R}^{N}$ & line-of-sight displacement at $N$ subsampled points \\
$\bfm \in \mathbb{R}^{P}$ & slip parameters (two components) plus nuisance terms \\
$\A \in \mathbb{R}^{N \times P}$ & design matrix: Green's functions, projected to LOS, plus ramp columns \\
$\Cd \in \mathbb{R}^{N \times N}$ & data covariance, block diagonal by track \\
$\Lop$ & discrete Laplace--Beltrami smoothing operator on the mesh \\
$\lambda$ & regularisation weight \\
$\M$ & regularised normal matrix, $\A'\Cd^{-1}\A + \lambda^2 \Lop'\Lop$ \\
$\R$ & model resolution matrix \\
$\hat{\bm{l}}$ & LOS unit vector, target to sensor \\
$L$ & element size, $\sqrt{2A}$ for a triangle of area $A$ \\
$r$ & three-dimensional source--receiver separation \\
\bottomrule
\end{tabular}
\caption{Notation.}
\label{tab:symbols}
\end{table}

\section{The forward problem}

\subsection{From interferometric phase to line-of-sight displacement}

An unwrapped interferogram carries phase $\phi$; the ground displacement
projected onto the radar line of sight is
\begin{equation}
d = -\frac{\lambda_{\text{radar}}}{4\pi}\,\phi ,
\label{eq:phase}
\end{equation}
positive \emph{toward} the sensor. The sign is not conventional decoration. With
interferograms formed as $\text{ref}\cdot\overline{\text{sec}}$ and ISCE3's
$\exp(-j4\pi r/\lambda)$ SLC phase convention, $\phi = (4\pi/\lambda)(r_{\text{sec}} -
r_{\text{ref}})$, so phase \emph{decreases} as the ground approaches the sensor
and the minus sign in \eqref{eq:phase} is required.

\note{This was wrong in every LOS product the package produced before 2026-07-28.}
The error was found not by a test --- the synthetic fixtures built the data with
$\G\bfm$ and inverted the same convention, so they were self-consistent under a
global flip --- but because inverting the known \emph{dextral} San Sebasti\'an
fault recovered \emph{sinistral} slip, with 86\% of elements pinned on the
polarity bound. It is now pinned by deriving the sign forward from a
displacement, and a mutation test confirms the pin bites: forcing the sign
positive drops the recovery test's variance reduction from $>95\%$ to 4.9\%.

\subsection{Green's functions}

Surface displacement from unit slip on element $j$, evaluated at station $i$, is
the integral of a point-source Green's function over the element,
\begin{equation}
\G_{ij} = \int_{T_j} \bm{g}(\bm{x}_i, \bm{\xi})\, N_j(\bm{\xi})\, \mathrm{d}A(\bm{\xi}),
\label{eq:greens}
\end{equation}
where $N_j$ is the basis function (constant on the element, or a piecewise-linear
``tent'' peaked at node $j$). The projection onto LOS gives one row of $\A$:
\begin{equation}
\A_{ij} = \hat{\bm{l}}_i \cdot \G_{ij}
        = l^E_i \G^E_{ij} + l^N_i \G^N_{ij} + l^U_i \G^U_{ij}.
\end{equation}

Two engines evaluate \eqref{eq:greens}. The \textbf{homogeneous half-space}
triangular dislocation element has a closed form and $\bm{g}$ never appears
explicitly. The \textbf{layered} engine tabulates $\bm{g}$ on a grid of
$(r, z)$ from EDGRN and evaluates \eqref{eq:greens} by quadrature:
\begin{equation}
\G_{ij} \;\approx\; A_j \sum_{q=1}^{n_q} w_q\, \bm{g}(\bm{x}_i, \bm{\xi}_q)\, N_j(\bm{\xi}_q).
\label{eq:quad}
\end{equation}

The layered engine matters because a homogeneous half-space gives the whole crust
one rigidity, while the shallowest few kilometres are much softer. Assuming
otherwise makes the same surface displacement require less shallow slip and more
deep slip --- a systematic bias in exactly the quantity a coseismic inversion is
for.

\subsection{Adaptive quadrature (novel)}
\label{sec:quad}

The reference implementations use a \emph{fixed} rule --- 91 points per
source--receiver pair in EDCMP's case. That includes this package's own direct
ancestor: \textbf{SlipSolve-Curve} (x3zou, MIT), the MATLAB code this work is a
port of and an extension to, itself descended from SlipSolve (ZeyuJin). It runs
the same nine-stage arrangement --- quadtree-sampled InSAR, a curved triangular
mesh, a fast homogeneous TDE pass and then a layered EDGRN/EDCMP inversion ---
at a fixed quadrature order throughout. That is wasteful by a wide
margin, because the integrand's difficulty depends entirely on how far the
observation is from the element \emph{in units of the element's own size}.

Measured on constant-slip triangular dislocations, the relative error of an
order-$k$ Gauss rule on \eqref{eq:quad} follows
\begin{equation}
\varepsilon \;\approx\; C \left(\frac{k\,r}{L}\right)^{-2},
\qquad C = 0.312 ,
\label{eq:err}
\end{equation}
with $L = \sqrt{2A_j}$ and $r$ the full three-dimensional separation from the
element centroid. The exponent is $-2$ (measured $-2.14$) because a triangle's
first moment vanishes about its own centroid, so the leading error term is second
order. Inverting \eqref{eq:err} for a requested tolerance $\tau$ gives the rule
actually used:
\begin{equation}
\boxed{\;
k_{ij} = \left\lceil \frac{L_j}{r_{ij}}\sqrt{\frac{\sigma C}{\tau}} \,\right\rceil,
\qquad \sigma = 4, \quad
k \in \left[1,\ \sqrt{N_{\max}}\right],\quad N_{\max}=900.
\;}
\label{eq:order}
\end{equation}
At the default $\tau = 3\times 10^{-3}$ this asks for $k \approx 20.4\,L/r$: order
one whenever the observation is more than about twenty element-sizes away, which
on a real problem is almost all of them.

Two details are load-bearing.

\paragraph{The safety factor $\sigma$.} The fit scatter around \eqref{eq:err} is
$1.89\times$ at one sigma, so a rule centred on the \emph{median} leaves half of
all pairs above the tolerance they asked for. Multiplying $C$ by $\sigma = 4$
raises the order by $\sqrt{\sigma}$ and takes the fraction of pairs exceeding
their tolerance from \textbf{51.7\% to 5.6\%}.

\paragraph{The near-field cap is a genuine, silent tolerance violation.} Inside
roughly $r \approx 2L$ the requested order exceeds $\sqrt{N_{\max}} = 30$ and is
refused. The smooth $1/(kr)$ law stops describing anything once the observation
is inside the element's own footprint, so the guarantee the method can honestly
make is \emph{never coarser than the fixed rule}, not \emph{always within
tolerance}. Capped pairs are now counted into the engine's provenance and warned
above 5\%; this is the method's one real weakness and it is reported rather than
hidden.

\paragraph{Performance.} On the six-scene problem (median $r/L = 38$), the
layered engine at $\tau = 3\times10^{-3}$ uses $2.59$ quadrature points per pair
over $3.4$ million pairs --- \textbf{48 times fewer} than the 91-point reference
rule --- and assembles in \textbf{0.86 times} the wall time of the closed-form
homogeneous engine. Against the reference rule it is $26\times$ faster, for a
slip difference of $0.47$~mm rms and a correlation of $0.9999997$.

\medskip
\noindent\textbf{Consequence.} There is no longer a computational reason to
prefer a homogeneous half-space. The bias it introduces is no longer being
bought with anything.

\section{Discretisation}

\subsection{Mesh and basis}

The fault is a triangulated surface built from a mapped surface trace and a dip
model. Slip has two components per basis function (strike-slip and dip-slip), so
$P = 2 n_{\text{basis}} + n_{\text{ramp}}$. Two bases are available:
\emph{element}-constant ($n_{\text{basis}} = n_{\text{elements}}$) and
\emph{node} ($n_{\text{basis}} = n_{\text{nodes}}$, piecewise-linear P1 tents).
Production uses the node basis; it gives a continuous slip field with roughly
half the parameters of the element basis on the same mesh.

\subsection{Regularisation}

Slip is regularised by a discrete Laplace--Beltrami operator on the mesh. For
interior edge $(i,j)$ the weight is the cotangent formula
\begin{equation}
w_{ij} = \tfrac{1}{2}\!\left(\cot\alpha_{ij} + \cot\beta_{ij}\right),
\end{equation}
$\alpha,\beta$ being the angles opposite the edge in the two adjacent triangles,
falling back to Floater's mean-value weights where the cotangent form is unsafe
(obtuse or boundary configurations). Rows are assembled as
$\Lop_{ii} = -\sum_{j} w_{ij}$, $\Lop_{ij} = w_{ij}$, then the whole operator is
normalised so that its largest absolute row sum equals~2 --- the scale of a plain
$[+1,-1]$ difference --- and repeated block-diagonally for the two slip
components.

\note{The units of the two weight families do not match, and mixing them was a
real bug.} Cotangents are dimensionless; mean-value weights carry $1/\text{length}$.
On a metre-scale mesh the ratio was measured at 2397, which is the median edge
length --- so a single node that fell back to the mean-value branch received a
smoothing row three orders of magnitude weaker than its neighbours and was
effectively unregularised. The symptom was one small patch of anomalous slip in
an otherwise clean model. The fix multiplies the mean-value fallback by the edge
length and a scale factor of $\tfrac12$. It required \emph{both} a fully starved
node \emph{and} a free-surface position to fire, which is why it survived so long.

\subsection{Nuisance parameters}

Each track carries an unknown reference level and a long-wavelength ramp
(orbital error, residual troposphere and ionosphere). These are solved
\emph{jointly} with slip rather than removed beforehand, as extra columns of
$\A$: a polynomial of chosen degree in the local $(x,y)$ of each track. The
production choice is quadratic, six columns per track, 36 columns over six
scenes.

Solving them jointly matters. Pre-removing a ramp by fitting it to the data
subtracts real deformation wherever the deformation is long-wavelength or reaches
the frame edge; solving jointly lets the data decide the split. It is not free
--- Section~\ref{sec:ramp} measures what the ramp absorbs even when solved
jointly --- but it is the lesser error.

\note{Do not merge scenes into one raster before inverting.} A mosaicked raster
has \emph{one} arbitrary constant and \emph{one} ramp, so a second scene's
independent nuisance is absorbed as fictitious slip. Measured on a controlled
synthetic with a planted right-lateral patch: two tracks concatenated as separate
observation sets recovered peak slip $1.587$~m against a truth of $1.626$~m
(correlation 0.979, moment ratio 1.001), while the same data mosaicked into one
raster returned $5.678$~m --- $3.5\times$ the truth (correlation 0.553, moment
ratio 1.671). \textbf{Variance reduction stayed at 98.4\%} for the wrong answer,
so the usual goodness-of-fit statistic does not catch this.

\section{The inverse problem}

\subsection{Data covariance}

InSAR noise is spatially correlated --- atmosphere, not shot noise --- so the
independent-errors assumption is wrong at exactly the scales slip inversion cares
about. We estimate an isotropic exponential covariance per track from the
far-field residual after removing a degree-2 trend:
\begin{equation}
\gamma(h) = n + s\left(1 - e^{-h/a}\right),
\qquad
C(h) = s\,e^{-h/a} + n\,\delta(h),
\label{eq:variogram}
\end{equation}
with sill $s$, range $a$ and nugget $n$ fitted to the empirical semivariogram.
The full $\Cd$ is block diagonal by track, since separate acquisitions have
independent atmospheres.

The generalised least-squares problem is then whitened by the Cholesky factor
$\Cd = \bm{F}\bm{F}'$, applied blockwise:
\begin{equation}
\tilde{\A} = \bm{F}^{-1}\A, \qquad \tilde{\bfd} = \bm{F}^{-1}\bfd ,
\end{equation}
by triangular solve rather than by forming $\Cd^{-1}$. Block diagonality is what
makes this affordable: six blocks of $\sim$1000--1600 at a few MB each, against
784--909~MB for a dense $\Cd$ at production size.

\subsection{The solve}

The estimate minimises
\begin{equation}
\boxed{\;
\hat{\bfm} = \arg\min_{\bfm \in \mathcal{B}}
\;\left\|\bm{F}^{-1}\left(\A\bfm - \bfd\right)\right\|_2^2
\;+\; \lambda^2 \left\|\Lop\bfm\right\|_2^2 \; }
\label{eq:objective}
\end{equation}
subject to box constraints $\mathcal{B}$: a \emph{polarity} constraint fixing the
sign of the strike-slip component to the known sense of the fault, and magnitude
bounds on each component. The bounded problem is solved by
\texttt{scipy.optimize.lsq\_linear}.

The unconstrained normal matrix,
\begin{equation}
\M = \A'\Cd^{-1}\A + \lambda^2\Lop'\Lop ,
\label{eq:normal}
\end{equation}
is what every resolution and uncertainty quantity below is built from. It is
Cholesky-factorised once per $\lambda$ and reused.

\note{The bounds interact with everything downstream and must be checked, not
assumed.} All the linear analysis in Sections~\ref{sec:res} and \ref{sec:unc}
describes the \emph{unbounded} problem. It is trustworthy only when few
parameters sit on a bound. On the current production model 6 of 374 nodes
(1.6\%) are active, all on the polarity constraint --- but an earlier run, before
the LOS sign fix, had \textbf{86\%} pinned, and there the linearisation would
have been meaningless. This is a property of the model, not of the code, and is
re-checked per run.

\section{Resolution analysis}
\label{sec:res}

\subsection{Model and data resolution}

With $\M$ from \eqref{eq:normal}, the model resolution matrix is
\begin{equation}
\R = \M^{-1}\A'\Cd^{-1}\A ,
\label{eq:R}
\end{equation}
so that $\hat{\bfm} = \R\,\bfm_{\text{true}}$ in the noiseless unbounded limit.
The data resolution (influence) matrix is
\begin{equation}
\bm{N} = \A\M^{-1}\A'\Cd^{-1},
\end{equation}
and the two satisfy
\begin{equation}
\operatorname{tr}(\R) = \operatorname{tr}(\bm{N}) \;\equiv\; \rho(\lambda),
\label{eq:trace}
\end{equation}
the \textbf{effective rank}: the number of parameters the data actually pays for.
Identity \eqref{eq:trace} is used as a runtime assertion --- it holds only if both
sides whiten the same $\A$, so a missed whitening site fails loudly instead of
silently biasing the analysis.

Effective rank, not $\lambda$, is the right axis for comparing configurations.
Two meshes at the same $\lambda$ are damped by different amounts; at the same
$\rho$ they are not. Every comparison in Section~\ref{sec:findings} is made at
matched $\rho$ for this reason.

\subsection{Averaging kernels}

Row $k$ of $\R$ is the averaging kernel for parameter $k$: it says which
neighbourhood of the true field the estimate at $k$ actually reports. Its width
is summarised by second moments of $|\R_{k,:}|$ in the two fault-surface
directions separately, giving anisotropic along-strike and down-dip resolving
lengths.

\note{These second moments are not Barnhart \& Lohman's metric, despite serving
the same purpose.} They fit a \emph{Gaussian} to row $k$ and take its
$1\sigma$. The averaging kernel here is strongly non-Gaussian: measured on the
six-scene model, \textbf{19.4\% of a row's $|\R|$ mass sits beyond $3\sigma$ of
the fitted Gaussian}, against the 0.27\% a Gaussian would put there --- a
$70\times$ excess. Second moments weight by distance squared, so they read those
tails and come out \textbf{3.09$\times$} the Gaussian $\sigma$ along strike
(39.3~km against 11.5~km). Trimming the weakest 10\% of each row brings the ratio
to 1.05, which identifies the tails as the entire difference.

\subsection{Resolution-driven meshing}

Two criteria are implemented for turning resolution into a target element size.

\paragraph{Diagonal (default).} Even out $\operatorname{diag}(\R)$ across the mesh:
\begin{equation}
h_{\text{new}} = h \left(\frac{\bar{r}}{r}\right)^{g},
\qquad g = 0.5,
\end{equation}
with $\bar{r}$ the median resolution, clipped to a factor of 3 per step and
smoothed over the mesh. Under-relaxing ($g<1$) is deliberate: resolution is itself
a function of the discretisation, so a full step chases its own tail.

\paragraph{Kernel width.} Set $h = f \times$ (kernel width), $f = 0.5$: two
elements per resolvable feature, the sampling analogue of Nyquist. This is the
closer analogue of the published method.

Element \emph{count} is held separately by an absolute integer budget, because
the size field sets the \emph{distribution} and not the scale. Targeting the
median looks self-normalising and is not: the resolution distribution is heavily
skewed, most elements hit the growth cap while few can shrink, and the mean
element size climbs every round. Measured without a budget, a loop coarsened 330
elements to 186, then 46, then 44, with variance reduction collapsing
$94.3\% \to 91.6\% \to \mathbf{0.74\%}$. It ran to completion, converged, and
returned nothing.

\subsection{The joint loop (novel)}

Barnhart \& Lohman (2010) size a mesh from resolution; Lohman \& Simons (2005)
subsample data from a predicted field. Their papers are explicit that the two
stay sequential and separate. The loop here drives \emph{both} from the same
$\R$, alternating:

\begin{enumerate}\itemsep2pt
\item Subsample each scene on the current model plus resolution field.
\item Assemble, solve, factorise $\M$ once.
\item Re-size elements from that $\R$; rebuild the mesh at fixed budget.
\item Repeat; hold the mesh fixed on the final round so the convergence
      measure is exact.
\end{enumerate}

Holding the mesh on the last round is necessary, not cosmetic. Re-meshing
perturbs a slip field by $0.1$--$1$~m at element boundaries whatever the model
does --- two models fitting the \emph{same} data on 240 and 232 elements agree to
a 2~mm median and still differ by $0.97$~m at the worst probe point --- so a
centimetre tolerance is unreachable while the mesh is moving.

\section{Uncertainty}
\label{sec:unc}

In a regularised inversion the error has two parts that trade off along the
L-curve, and reporting only one is misleading.

\subsection{Variance}

Propagating data noise through the linearised estimator gives
\begin{equation}
\Cm = \M^{-1}\A'\Cd^{-1}\A\,\M^{-1} .
\label{eq:post}
\end{equation}
When the estimator's weighting does not match the true noise --- which is the
production case, since the shipped model is solved with plain weights while the
noise is correlated --- the correct expression is the \emph{sandwich}
\begin{equation}
\boxed{\;\Cm = \M^{-1}\,\A'\W\Cd\W\,\A\,\M^{-1}\;}
\label{eq:sandwich}
\end{equation}
with $\W$ the diagonal weight matrix the solve actually used. Note $\W$ appears
\emph{twice} on each side; applying it once understates the result by a
constant factor (measured: $28.78\times$ too large in the first implementation,
with the correlation against Monte Carlo unchanged at 0.978, which is what
identified it as a pure scale error rather than a structural one).

\subsection{Bias}

The regulariser's pull on the model is
\begin{equation}
\bm{b} = (\R - \bm{I})\,\bfm ,
\label{eq:bias}
\end{equation}
and $1 - \R_{kk}$ is directly the fraction of the estimate at $k$ that came from
smoothing rather than from data.

\note{Equation \eqref{eq:bias} must be evaluated at a \emph{fixed} reference
model, not at the recovered model.} Evaluated at $\hat{\bfm}(\lambda)$ the bias
appears to \emph{fall} with $\lambda$ (measured $0.317 \to 0.098$) simply because
the recovered model shrinks; and even the relative bias is non-monotonic
($0.203 \to 0.431 \to 0.290$) because $\hat{\bfm}$ rotates toward the directions
$\R$ preserves. Only against a fixed reference does the curve mean what it
appears to mean.

\subsection{Monte Carlo}

The honest check draws realisations of correlated noise from the fitted $\Cd$,
adds them to the data, and re-solves through the \emph{actual bounded} solver,
which captures the bounds and the nonlinearity the linearisation discards. The
Green's matrix is reused across realisations, so this is $N$ solves and not
$N$ assemblies.

\note{Not bootstrap.} A naive observation resample assumes independence, which
the measured 22--33~km correlation length says is false, and a block bootstrap
would need blocks larger than that --- leaving too few independent blocks in a
5700-sample set. Drawing from a covariance we actually measured is strictly
better here.

\section{Findings}
\label{sec:findings}

\subsection{The noise structure of the six scenes}

Table~\ref{tab:cov} gives the fitted exponential covariances. The striking
result is the third column.

\begin{table}[htbp]
\centering
\small
\begin{tabular}{lrrr}
\toprule
scene & sill (mm) & nugget (mm) & range (km) \\
\midrule
D135 & 13.4 & 0.42 & 31.3 \\
D134 & 16.3 & 0.52 & 32.6 \\
A061 & 18.1 & 1.19 & 22.5 \\
A162 & 20.7 & 2.25 & 22.7 \\
D126 & 27.3 & 3.27 & 23.2 \\
D054 & 40.9 & 1.29 & 29.3 \\
\bottomrule
\end{tabular}
\caption{Exponential covariance fitted to each scene's far-field residual after
a degree-2 detrend. Sill and nugget are quoted as standard deviations.}
\label{tab:cov}
\end{table}

The sill varies by a factor of three across scenes and the nugget by nearly
eight, but the \textbf{correlation range is 22--33~km on all six} --- across two
different processing chains (NISAR GSLC and ALOS). That is a property of the
atmosphere over this region, not of the processing, and it is the number that
makes a block bootstrap infeasible.

\subsection{The data covariance does not improve the slip estimate}

This was the load-bearing gap identified before the work started, and closing it
produced a negative result. Against a synthetic with known truth, over five
seeds, the whitened solve recovered slip at \textbf{0.147~m rms error} against
\textbf{0.069~m} for the plain-weighted solve --- worse, consistently, 5 seeds
out of 5.

The reason is not a bug. Generalised least squares is minimum-variance only among
\emph{unbiased} estimators. This estimator is neither: it is bound-constrained
and, at the operating point, roughly 97\% regulariser. Whitening correctly
down-weights spatially redundant information, but what it hands back to the
regulariser is a weaker data constraint, and the regulariser is the larger term.

What $\Cd$ \emph{is} good for is honest error bars --- the sandwich
\eqref{eq:sandwich} needs it --- and a $\chi^2$ per degree of freedom that means
something. It is retained for those, not for the point estimate.

\subsection{The element budget and \texorpdfstring{$\lambda$}{lambda}}

The inherited budget (1696 elements, 1836 slip parameters) traced back to a
default edge length, not to a measurement. Against it, $\operatorname{tr}(\R) =
38.7$: the data paid for about 1 parameter in 47.

Re-measured, the L-curve corner sits at effective rank $\approx 32$ on the
1696-element mesh and $\approx 21$ on the 651-element unstructured mesh. The
budget was \textbf{11--57 times larger than the data supports}. Production now
runs 660 elements (748 parameters), which reaches rank 41.5 at $\lambda = 30$ and
so clears the corner with headroom, while 416 elements caps at 25.7 and cannot.

\subsection{The ramp absorbs signal, and how much depends on \texorpdfstring{$\lambda$}{lambda}}
\label{sec:ramp}

Ramp degree buys fit monotonically (variance reduction $96.50 \to 99.55 \to 99.66
\to 99.71$ for offset / linear / quadratic / cubic) while moment changes by
$1.00 \to 1.56 \to 1.65 \to 1.48$ relative to the offset-only case. Fit alone
cannot choose the degree.

The synthetic recovery experiment measures directly what the ramp eats, because
truth is known and the data are noiseless: any ramp coefficient returned is
signal the ramp has absorbed. On the 40~km magnitude checkerboard the maximum
absorbed amplitude falls by a factor of ten as damping is relaxed
(Table~\ref{tab:ladder}). \textbf{Ramp leakage is largely a symptom of
over-damping}: when the slip parameterisation cannot fit the data, the ramp
absorbs the misfit.

\begin{table}[htbp]
\centering
\small
\begin{tabular}{rrrrrrrr}
\toprule
& & \multicolumn{3}{c}{pattern correlation} & \multicolumn{3}{c}{max $|$ramp$|$ (m)} \\
\cmidrule(lr){3-5}\cmidrule(lr){6-8}
rank & $\lambda$ & 20 km & 40 km & 60 km & 20 km & 40 km & 60 km \\
\midrule
 9.7 & 388 & 0.018 & 0.278 & 0.250 & 0.730 & 0.670 & 0.728 \\
20.6 & 144 & 0.033 & 0.437 & 0.423 & 0.471 & 0.392 & 0.443 \\
39.8 &  50 & 0.070 & 0.615 & 0.634 & 0.250 & 0.167 & 0.195 \\
79.1 &  13 & 0.168 & 0.726 & 0.751 & 0.118 & 0.070 & 0.079 \\
\bottomrule
\end{tabular}
\caption{Recovery of magnitude checkerboards against effective rank, on noiseless
synthetic data with real coverage and the unstructured layout. Correlation is
computed on the oscillating part, after removing each field's own mean.
Production operates at rank $\approx 21$, the second row. $\lambda$ is quoted for
the 40~km case; the other two differ by under 1\%.}
\label{tab:ladder}
\end{table}

Three things fall out.

\paragraph{There is a resolution cliff between 20 and 40~km.} The 20~km
checkerboard \textbf{never recovers}: correlation runs from 0.018 to 0.168 across
an eight-fold change in effective rank, and even at rank 79 the recovered field
(Figure~\ref{fig:ladder}, top row) bears no resemblance to the target. Relaxing
$\lambda$ does not buy it back, so this is a genuine information limit and not a
regularisation choice.

\paragraph{Above that cliff, wavelength stops mattering and damping takes over.}
The 40 and 60~km cases track each other to within 0.02 in correlation at every
rung --- 0.28/0.25, 0.44/0.42, 0.62/0.63, 0.73/0.75 --- while each varies by a
factor of 2.6 across the ladder. Between 40 and 60~km the answer is set by
$\lambda$, not by the target.

\paragraph{The kernel width marks ``half-recovered'', not ``recovered''.} At the
production operating point the measured along-strike kernel width is $\approx
40$~km, and a 40~km checkerboard comes back at correlation 0.44. So the width is
a fair marker of where recovery becomes partial, and an optimistic one if read as
where recovery becomes clean.

\medskip
\noindent
A caution on metrics: at 20~km the \emph{amplitude ratio} sits at 0.88 while the
correlation is 0.03. The model carries very nearly the right amount of variance
and puts almost none of it in the right place. Amplitude ratio alone is not a
recovery statistic.

\begin{figure}[htbp]
\centering
\includegraphics[width=\textwidth]{../venezuela_six_scene_run/results_recovery/resolution_ladder.png}
\caption{The resolution ladder. Rows are planted wavelength, columns are
effective rank; the leftmost column is truth. Same noiseless data, same
estimator, same mesh budget --- only $\lambda$ varies along a row. The 20~km
target never resolves; the 40 and 60~km targets are limited by damping rather
than by wavelength. Each panel has its row's colour scale, with the uniform
dextral base removed.}
\label{fig:ladder}
\end{figure}

\subsection{Adaptive quadrature}

Reported in Section~\ref{sec:quad}: 48$\times$ fewer quadrature points than the
reference rule, layered assembly at 0.86$\times$ the cost of the closed-form
homogeneous engine, median relative error $7.2\times10^{-4}$ against real
Crust2.0 tables. This is the cleanest positive result in the study.

\subsection{The joint loop is an efficiency method, not an accuracy method}

Four arms at one budget, one round count, compared at matched effective rank:
a fixed uniform mesh (null), resolution-sized once then sampled separately (the
published arrangement), the joint loop, and the joint loop with the kernel-width
metric. The joint loop gains \textbf{3.6\% over the null} at the corner rank ---
and costs \textbf{32\% more}.

The earlier belief that it was also cheaper was an artefact of a defect: the loop
was performing seven Cholesky factorisations of a $\sim$1870$^2$ matrix per round
where one suffices. With that fixed, the cost advantage disappears.

The kernel-metric arm \emph{lost} at every effective rank, including to a fixed
mesh at low rank. Section~\ref{sec:res} explains why, and the explanation
exonerates the published method: our kernel width is a second moment, not their
Gaussian fit, and the reduction \texttt{min} selects the down-dip width on 100\%
of parameters, so the along-strike number never reaches the size field at all.
This is evidence against \emph{this estimator of the kernel}, not against
Barnhart \& Lohman's method, which has not been run here.

\subsection{Uncertainty on the production model}

The analytic sandwich \eqref{eq:sandwich} reproduces a 40-draw Monte Carlo
through the bounded solver at a median ratio of \textbf{1.15} and a correlation
of \textbf{0.79} over 374 nodes, in under a second against roughly 40 minutes.
The shipped model has peak slip $\mathbf{2.28 \pm 0.07}$~m, with a mean
$\sigma$ of 0.035~m analytically and 0.028~m by Monte Carlo.

\note{These are variance only, and variance is the smaller half.} The
regularisation bias \eqref{eq:bias} exceeds the propagated data variance by an
order of magnitude at every usable $\lambda$. An error bar quoting only $\sigma$
would be badly over-confident. The correct presentation is the bias--variance
trade-off against $\lambda$, not a single number: moving $\lambda$ moves both, so
they cannot be added into one interval.

\subsection{Synthetic recovery, and a test that measured its own bound}
\label{sec:recovery}

Truth was planted on the mesh, forward-modelled through the real look geometry
and coverage masks, and inverted with the production configuration. Noiseless by
design: a checkerboard's adjacent lobes cancel at the surface, leaving only
10--18~mm rms of signal against a fitted atmosphere of 12--41~mm, so noise would
swamp it and the test would report ``nothing recovers''.

\note{The first suite was invalid, and every summary statistic stayed plausible
while it was.} The textbook checkerboard alternates sign; production runs a
polarity constraint forcing strike-slip $\le 0$. \textbf{43.8\% of the planted
model was in a region the solver is forbidden to enter.} Every recovered model
came back single-signed at all 18{,}000 probe points, with 25--32\% sitting at
exactly zero (Figure~\ref{fig:recovery}). The scores were measuring the bound.

The fix keeps the estimator and changes the target: alternate the \emph{magnitude}
of dextral slip rather than its sign,
\begin{equation}
s(\bm{x}) = -\left[\,b + c\,\sin(\pi x_s/\ell)\sin(\pi x_z/\ell)\,\right],
\qquad b = 1.25,\; c = 1.0,
\end{equation}
spanning $-0.25$ to $-2.25$~m --- a factor of nine between weak and strong cells,
entirely inside the bound. Sixteen inversions were run on it: three cell sizes,
two layouts at 40~km, each at four points on the rank ladder
(Table~\ref{tab:ladder}, Figure~\ref{fig:ladder}).

Two conclusions survive the confound, because the bound applied identically to
every arm:
\begin{itemize}\itemsep2pt
\item \textbf{Mesh layout is not the lever, at any damping.} Row-structured
      and unstructured layouts on the 40~km target agree to within 3\% at every
      rung of the ladder --- correlation 0.278/0.278, 0.422/0.437, 0.610/0.615,
      0.729/0.726 --- and the sign of the difference \emph{flips} at the last
      one, so it is scatter and not a trend. On real data, compared only to each
      other, the same two layouts differ by 0.78~m locally.
\item \textbf{The quadratic ramp absorbed 0.38--0.39~m} of a $\pm1.94$~m
      checkerboard under real coverage, and 0.14--0.16~m under full coverage.
      Coverage gaps roughly double what the ramp takes.
\end{itemize}

A separate check confirms the mesh is not the limiting factor: projected onto the
660-element mesh's nodes, the planted checkerboards retain 0.993, 0.999 and 0.993
of their amplitude at 20, 40 and 60~km cell size. Any amplitude loss in the
recovery is the regulariser, not the discretisation.

\begin{figure}[htbp]
\centering
\includegraphics[width=\textwidth]{../venezuela_six_scene_run/results_recovery/recovery_checker40km.png}
\caption{Planted sign-alternating checkerboard (top) against four recovered
models. Left column shares the truth's colour scale; right column is
autoscaled per panel. Every recovered model is single-signed --- the polarity
constraint forbids the positive lobes, so the positive cells are returned as
exactly zero. This is the figure that identified the test as invalid.}
\label{fig:recovery}
\end{figure}

\subsection{The L-curve corner is not where the error is smallest}
\label{sec:optimum}

The ladder above is noiseless, so it measures only the \emph{bias} half of the
error. That is exactly the half the L-curve corner is not trying to minimise, so
the two must be compared under noise before either can be believed. Five
independent realisations were drawn from the covariances of
Table~\ref{tab:cov}, added to the 40~km target, and inverted over six rungs
(Table~\ref{tab:optimum}, Figure~\ref{fig:optimum}).

\begin{table}[htbp]
\centering
\small
\begin{tabular}{rrrrrrr}
\toprule
rank & $\lambda$ & correlation & rel.\ error & amplitude & max $|$ramp$|$ (m) \\
\midrule
 10 & 381.0 & $0.275 \pm 0.014$ & $0.618 \pm 0.002$ & 0.904 & 0.667 \\
 21 & 142.5 & $0.438 \pm 0.024$ & $0.494 \pm 0.005$ & 0.950 & 0.411 \\
 39 &  51.4 & $0.619 \pm 0.019$ & $0.382 \pm 0.007$ & 1.082 & 0.290 \\
\textbf{79} & \textbf{12.8} & $\bm{0.715 \pm 0.022}$ & $\bm{0.353 \pm 0.015}$ & 1.287 & 0.263 \\
129 &   3.8 & $0.692 \pm 0.016$ & $0.392 \pm 0.025$ & 1.416 & 0.269 \\
198 &   1.1 & $0.519 \pm 0.047$ & $0.609 \pm 0.043$ & 1.878 & 0.271 \\
\bottomrule
\end{tabular}
\caption{Five noise realisations of the 40~km magnitude checkerboard, real
coverage, unstructured layout. Uncertainties are the seed spread. Every solve
converged.}
\label{tab:optimum}
\end{table}

\noindent
\textbf{The error curve turns over at effective rank 79 --- 3.8 times the
L-curve corner --- and the corner is 40\% worse.} The seed spread is 0.015 at the
optimum against a 0.14 gap to the corner, so this is a property of the problem
and not of one realisation.

The three panels of Figure~\ref{fig:optimum} separate the two error terms
cleanly. Pattern correlation peaks at the same rank the total error minimises.
Amplitude ratio, by contrast, climbs \emph{monotonically} through 1.0 and reaches
1.88 at rank 198 --- that is the variance term, and it is why amplitude ratio
must never be read as a recovery statistic on its own.

\note{This does not make the L-curve wrong; it makes it a different criterion.}
The L-curve reads the bend in misfit against roughness and needs no truth, which
is exactly why it is used. What this measures is that on this problem the bend
sits at about a quarter of the rank the data can actually carry. Two caveats
bound the claim: it is measured on a 40~km checkerboard, and a target with
different spectral content may optimise elsewhere; and the ramp is solved
jointly, which absorbs much of the correlated noise (max $|$ramp$|$ is roughly
flat at 0.26--0.29~m from rank 39 upward, against 0.07~m with no noise) and so
flattens the variance penalty that would otherwise punish low damping sooner.

\begin{figure}[htbp]
\centering
\includegraphics[width=\textwidth]{../venezuela_six_scene_run/results_recovery/lambda_optimum.png}
\caption{Choosing $\lambda$ against known truth. Left: total error against truth,
minimised at rank 79. Centre: pattern correlation, peaking in the same place.
Right: amplitude ratio, climbing monotonically past 1.0 --- the variance term
made visible. Error bars are the spread over five noise realisations.}
\label{fig:optimum}
\end{figure}

\subsection{Two things that do not help: more elements, and more data}
\label{sec:nohelp}

The 20~km target fails at every damping. Two candidate explanations were tested
directly and both are refuted, which leaves an information limit.

\paragraph{Doubling the element budget changes nothing.} At 660 elements the
correlations across the ladder are 0.018 / 0.034 / 0.070 / 0.169; at 1320 they
are 0.018 / 0.035 / 0.064 / 0.156 --- indistinguishable, and marginally worse
where it matters. Consistent with the discretisation ceiling, which already put
99.3\% of that target's amplitude inside the coarser mesh: the extra parameters
land in the null space, because effective rank is set by the data.

\paragraph{Full coverage does not help either, and at 40~km it hurts.} Removing
every coverage mask raises the observation count from about 7\,200 to about
17\,200 and moves the 20~km correlations by less than 0.002 at the first three
rungs. At 40~km the effect is \emph{negative} and reproduces on both layouts:

\begin{center}
\small
\begin{tabular}{lrrrr}
\toprule
full minus real coverage & rank 10 & rank 21 & rank 40 & rank 79 \\
\midrule
40 km, unstructured & $+0.003$ & $-0.043$ & $-0.053$ & $-0.019$ \\
40 km, rows         & $-0.009$ & $-0.040$ & $-0.046$ & $-0.015$ \\
20 km, unstructured & $+0.000$ & $+0.002$ & $+0.001$ & $-0.021$ \\
\bottomrule
\end{tabular}
\end{center}

\noindent
The mechanism is the comparison itself: at \emph{matched effective rank} the same
rank budget is spread over $2.4\times$ as many observations, so each is damped
harder. The operational consequence is worth stating plainly --- \textbf{more
data does not buy resolution unless $\lambda$ is relaxed to spend it}, which is
the same lesson as Section~\ref{sec:optimum} arriving from the other direction.

\subsection{Summary: \texorpdfstring{$\lambda$}{lambda} dominates}

Four independent lines of evidence say the regularisation weight determines the
answer and the other design choices are second-order:

\begin{enumerate}\itemsep2pt
\item Four mesh resolutions at matched effective rank agree on rms misfit to
      0.3\% and on $M_w$ to 0.005.
\item Four mesh designs on the real six-scene data gave variance reduction
      99.65--99.67\% and $M_w$ 7.594--7.607.
\item Two layouts, scored against known truth, agree to 1\%.
\item Recovery of a fixed target varies from correlation 0.28 to 0.73 across the
      rank ladder alone, while swapping the mesh layout at fixed rank moves it by
      0.015 and changing the target wavelength from 40 to 60~km moves it by 0.02
      (Table~\ref{tab:ladder}).
\end{enumerate}

\noindent
The fourth is the sharpest: on the same data, $\lambda$ moves the answer roughly
\textbf{thirty times} as far as the mesh layout does. And
Section~\ref{sec:optimum} turns that from a ranking into a number ---
the operating point everything else was measured at was itself
\textbf{40\% worse than the best available}, and no mesh, layout, budget or
coverage change recovers anything like that much.

Ortega-Culaciati et al.\ (2021) reach the same conclusion from theory and respond
by fixing a fine mesh and varying the prior instead. Our data supports that
choice. The honest headline is that the joint loop is an \emph{efficiency}
method, that the adaptive quadrature is a genuine and substantial computational
result, and that the model itself is set by the prior.

\section{Bugs worth recording}

Each of these produced plausible output while being wrong, which is the only
reason they are worth listing.

\begin{enumerate}\itemsep3pt
\item \textbf{LOS sign inverted} (Section~2.1). Every LOS product mirrored. Found
      by a dextral fault recovering sinistral slip, not by a test; the tests were
      self-consistent under a global flip.
\item \textbf{Smoothing operator units} (Section~3.2). Cotangent weights mixed
      with mean-value weights, a factor of the median edge length (2397) apart.
      Symptom: one patch of anomalous slip.
\item \textbf{Sandwich covariance} (Section~6.1). Weight matrix applied once
      instead of twice; $28.78\times$ too large, correlation against Monte Carlo
      unchanged.
\item \textbf{Bias evaluated at a moving reference} (Section~6.2). Produced a
      bias curve with the wrong sign of slope.
\item \textbf{Checkerboard against a polarity bound} (Section~\ref{sec:recovery}).
      An entire recovery suite measuring its own constraint.
\item \textbf{Polarity diagnostic inverted.} The check written to catch (5)
      tested the reachable set instead of the forbidden one, reporting ``100\%
      forbidden'' for a target lying entirely inside the bound. Diagnostic only;
      no score affected.
\item \textbf{Coverage mask reached the forward model through the look vectors,}
      not through the LOS field --- so a ``full coverage'' synthetic world was
      byte-identical to the masked one, both at 24\% coverage, until the look
      vectors were filled too.
\end{enumerate}

\section{Open items}

\begin{itemize}\itemsep2pt
\item \textbf{Gaussian-fit kernel width.} The one variant of the published metric
      that has never been implemented, and the one the diagnosis in
      Section~\ref{sec:res} points at.
\item \textbf{The near-field quadrature cap} (Section~\ref{sec:quad}) is a known
      tolerance violation inside $r \approx 2L$. Bounded and reported, not fixed.
\item \textbf{Multi-fault geometry.} Blocked on the boundary-node mask, which
      tests global extrema, so a second fault's ends silently lose their
      zero-slip rows. The fix is a per-component label; everything downstream
      falls out unchanged.
\item \textbf{The A/B loop comparison predates the smoothing-operator fix} and
      has not been re-run.
\item \textbf{Effect of the mesh on where slip goes at depth.} The layouts agree
      on scalars and on truth-referenced error, but 64\% of their mutual
      disagreement sits below 20~km depth, where $\operatorname{diag}(\R)\approx 0$.
      Neither can be called right there; the honest statement is that the data
      does not constrain it.
\end{itemize}

\section*{References}
\small
\begin{itemize}\itemsep1pt
\item Barnhart, W.\ D.\ \& Lohman, R.\ B.\ (2010). Automated fault model discretization for
      inversions for coseismic slip distributions. \emph{JGR Solid Earth} 115, B10419.
\item Lohman, R.\ B.\ \& Simons, M.\ (2005). Some thoughts on the use of InSAR data to
      constrain models of surface deformation. \emph{G-cubed} 6, Q01007.
\item Nikkhoo, M.\ \& Walter, T.\ R.\ (2015). Triangular dislocation: an analytical,
      artefact-free solution. \emph{Geophys.\ J.\ Int.} 201, 1119--1141.
\item Ortega-Culaciati, F.\ et al.\ (2021). Asperities and barriers on the seismogenic
      zone: a Bayesian perspective. \emph{JGR Solid Earth}.
\item Wang, R., Martin, F.\ L.\ \& Roth, F.\ (2003). Computation of deformation induced by
      earthquakes in a multi-layered elastic crust (EDGRN/EDCMP).
      \emph{Computers \& Geosciences} 29, 195--207.
\item x3zou. \emph{SlipSolve-Curve}: reproducible MATLAB coseismic-slip inversion
      on triangular fault meshes (MIT; after SlipSolve by ZeyuJin).
      \url{https://github.com/x3zou/SlipSolve-Curve} --- \textbf{the direct
      ancestor of the code described here}.
\end{itemize}

\end{document}
